%--------------------------------------------------------------------
\subsection{Collinear Spin Extension}
\label{sec:collinear-spin-extension}
%--------------------------------------------------------------------

The spin-polarized extension used in the present work is collinear.
A fixed spin quantization axis is chosen, and the spin index
$\sigma\in\{\uparrow,\downarrow\}$ labels two spin channels.  For each
$\mathbf k$ point, the one-particle density matrix and Hamiltonian have
a block-diagonal spin structure,
\begin{equation}
  \mathbf P_{\mathbf k}(t)
  =
  \begin{pmatrix}
    \mathbf P_{\mathbf k}^{\uparrow}(t) & 0 \\
    0 & \mathbf P_{\mathbf k}^{\downarrow}(t)
  \end{pmatrix},
  \qquad
  \mathbf H_{\mathbf k}(t)
  =
  \begin{pmatrix}
    \mathbf H_{\mathbf k}^{\uparrow}(t) & 0 \\
    0 & \mathbf H_{\mathbf k}^{\downarrow}(t)
  \end{pmatrix}.
  \label{eq:collinear_spin_block_structure}
\end{equation}
Thus there are no off-diagonal spinor blocks in the present formulation.
The two spin channels are propagated as separate density matrices, each
obeying the Liouville-von Neumann equation introduced in
Sec.~\ref{sec:density-matrix-propagation},
\begin{equation}
  \imath\hbar
  \frac{d\mathbf P_{\mathbf k}^{\sigma}(t)}{dt}
  =
  \left[
    \mathbf H_{\mathbf k}^{\sigma}(t),
    \mathbf P_{\mathbf k}^{\sigma}(t)
  \right].
  \label{eq:lvn_collinear_spin}
\end{equation}

The spin-resolved electron densities are reconstructed from
$\mathbf P_{\mathbf k}^{\sigma}(t)$ as described in
Sec.~\ref{sec:density-matrix-propagation}.  The total electron density is
\begin{equation}
  \rho(\mathbf r,t)
  =
  \rho^{\uparrow}(\mathbf r,t)
  +
  \rho^{\downarrow}(\mathbf r,t).
  \label{eq:total_density_spin}
\end{equation}

In the collinear Kohn--Sham Hamiltonian, the local effective potential
for spin channel $\sigma$ can be written schematically as
\begin{equation}
  \hat V_{\rm loc}^{\sigma}
  [\rho^\uparrow,\rho^\downarrow](t)
  =
  \hat V_{\rm ion}^{\rm loc}
  +
  \hat V_{\rm H}[\rho](t)
  +
  \hat V_{\rm xc}^{\sigma}
  [\rho^\uparrow,\rho^\downarrow](t).
  \label{eq:vloc_collinear_spin}
\end{equation}
The Hartree potential depends on the total electron density
$\rho=\rho^\uparrow+\rho^\downarrow$, while the exchange-correlation
potential is spin-channel dependent and is a functional of both spin
densities.  In the absence of spin-orbit coupling, the nonlocal
pseudopotential terms considered here are diagonal in the collinear spin
basis and are treated as described in
Secs.~\ref{sec:periodic-velocity-gauge} and~\ref{sec:current-nlpp}.

Equation~\eqref{eq:lvn_collinear_spin} therefore represents separate
propagation within each spin channel, but the channels are not
physically independent.  They are coupled self-consistently through the
density-dependent potentials:
\begin{equation}
  \left\{
    \mathbf P_{\mathbf k}^{\uparrow},
    \mathbf P_{\mathbf k}^{\downarrow}
  \right\}
  \rightarrow
  \left\{
    \rho^\uparrow,
    \rho^\downarrow
  \right\}
  \rightarrow
  \hat V_{\rm H}[\rho^\uparrow+\rho^\downarrow],
  \hat V_{\rm xc}^{\sigma}[\rho^\uparrow,\rho^\downarrow]
  \rightarrow
  \left\{
    \mathbf H_{\mathbf k}^{\uparrow},
    \mathbf H_{\mathbf k}^{\downarrow}
  \right\}.
  \label{eq:collinear_spin_coupling_cycle}
\end{equation}

The same spin decomposition applies to the macroscopic charge current.
For each spin channel,
\begin{equation}
  J_\alpha^{\sigma}(t)
  =
  \frac{1}{V_{\rm cell}}
  \sum_{\mathbf k}
  w_{\mathbf k}\,
  {\rm Tr}\!
  \left[
    \mathbf P_{\mathbf k}^{\sigma}(t)
    \hat I_{\mathbf k,\alpha}^{\sigma}(t)
  \right],
  \label{eq:spin_resolved_current}
\end{equation}
where the charge-current operator is the spin-diagonal operator defined
in Sec.~\ref{sec:current-nlpp}.  The total macroscopic charge-current
density is the sum over the two collinear spin channels,
\begin{equation}
  J_\alpha(t)
  =
  \sum_{\sigma=\uparrow,\downarrow}
  J_\alpha^{\sigma}(t).
  \label{eq:total_current_spin_sum}
\end{equation}
The individual $J_\alpha^{\sigma}(t)$ may also be analyzed separately
to obtain spin-resolved optical response within the fixed collinear
basis.

The present formulation does not include off-diagonal spin coherences,
transverse spin dynamics, general two-component spinors, noncollinear
magnetization dynamics, or spin-orbit-induced spin mixing.  Such
extensions require a spinor Hamiltonian with off-diagonal spin blocks
and lie outside the collinear spin-polarized RT-TDDFT framework used in
this manuscript.  This spin-resolved collinear formulation provides the
theoretical basis for analyzing the spin-up and spin-down optical
channels of the NV$^{-}$ center discussed below.
